\section{六、灵敏度分析与稳健性检验}\label{ux516dux7075ux654fux5ea6ux5206ux6790ux4e0eux7a33ux5065ux6027ux68c0ux9a8c}

\subsection{\texorpdfstring{6.1 安全库存调节系数 \(z\)
的全局灵敏度分析}{6.1 安全库存调节系数 z 的全局灵敏度分析}}\label{ux5b89ux5168ux5e93ux5b58ux8c03ux8282ux7cfbux6570-z-ux7684ux5168ux5c40ux7075ux654fux5ea6ux5206ux6790}

安全库存系数 \(z\) 是权衡库存持有水平与缺货风险的核心决策参数。将 \(z\)
在区间 \([0.0, 3.0]\) 内以步长 \(0.1\) 进行全局连续扰动扫描，考察其对 6
种物料平均服务水平与平均库存量的边际影响。

如图 6.1(a) 所示： 1. 当 \(z\) 从 \(0.0\) 增至 \(1.2\)
时，平均服务水平由 \(88.5\%\) 陡增至
\(96.8\%\)，表现出极高的控制敏感度； 2. 当 \(z > 1.8\)
之后，服务水平逐渐饱和收敛于 \(99.1\%\)
附近，而平均在库库存量则呈线性单调攀升（从 11.2 件增至 48.6
件），边际保障效益显著递减。这表明推荐选取的 \(z \in [0.8, 1.4]\)
处于极佳的边际效益黄金区间。

\subsection{6.2 突发需求脉冲冲击与系统稳健性（Stress
Testing）}\label{ux7a81ux53d1ux9700ux6c42ux8109ux51b2ux51b2ux51fbux4e0eux7cfbux7edfux7a33ux5065ux6027stress-testing}

为检验排产系统在面对非平稳极端异常扰动时的鲁棒性，在第 120 周与第 121
周对 6 种物料人为注入 3 倍于实际需求的突发脉冲冲击（Demand Surge
Shock），观察系统的抗毁度与动态收敛恢复时间（Settling Time）。

如图 6.1(b) 所示： 1. 在第 120～121
周脉冲冲击期间，由于生产提前期时滞，当周服务水平出现结构性下挫； 2.
冲击结束后，得益于动态 Base-Stock 的负反馈差值补偿，排产系统在 2～3
周内迅速完成补库拉升，第 122～177 周的平均服务水平恢复至
\(94.6\%\)～\(100.0\%\)，全期综合服务水平依然稳定在
\(89.8\%\)～\(100.0\%\)（全部远超 \(85\%\)
红线）。实测证明该排产算法具备极强的抗扰动自愈能力。

\begin{figure}[H]
\centering
\includegraphics[width=0.88\textwidth]{figures/fig6_sensitivity_and_robustness.pdf}
\caption{安全库存系数全局灵敏度分析 (a) 与突发需求冲击稳健性恢复检验 (b)}
\label{fig:fig6}
\end{figure}
